\chapter{Maximum Likelihood Estimation}
\label{ch:mle}

The estimators of the previous chapter were built one at a time, by ad hoc
reasoning. We now meet a single, general recipe for producing good
estimators—\emph{maximum likelihood}. The principle is disarmingly simple:
choose the parameter value that makes the observed data most probable.

\section{The Maximum Likelihood Principle}

\begin{principle}[Maximum likelihood]
Given a dataset, estimate the unknown parameter by the value that makes the
observed data as likely as possible.
\end{principle}

A small story fixes the idea. A dealer is offered one of two batches of chips:
one is about $10\%$ defective, the other about $50\%$, but the seller has lost
track of which is which. The dealer tests $10$ chips from a batch and finds
exactly one defective. Which batch is it more likely to be?

\begin{example}[Which batch?]
Test $10$ chips; let $R_i=1$ if chip $i$ is defective. With one defective among
ten, the probability of the observed pattern is $p(1-p)^9$. Compare the two
candidate values of $p$:
\[
  \text{($10\%$ batch)}\ \ (0.1)(0.9)^9\approx 0.039,\qquad
  \text{($50\%$ batch)}\ \ (0.5)(0.5)^9\approx 0.00098.
\]
The observed data are about forty times more likely under the $10\%$ batch, so
the dealer should believe it is that one. This is maximum likelihood reasoning in
miniature.

If instead the \emph{first three} chips were defective, the data probability
becomes $p^3(1-p)^7$, now favoring the $50\%$ batch ($\approx 0.00098$ versus
$\approx 0.00048$): three early defectives are better explained by the
high-defect batch. The data, not a prior hunch, decide.
\end{example}

\begin{example}[Which die?]
A die $D_1$ has five red faces and one white; $D_2$ has one red and five white.
One die is chosen and rolled until red appears, three separate times, giving
first-red on the $3$rd, $5$th, and $4$th rolls. Which die is it likely to be?

\solution The rolls are independent $\Geo(p)$ with $p$ the chance of red, so the
likelihood is
\[
  L(p)=(1-p)^2 p\cdot(1-p)^4 p\cdot(1-p)^3 p=p^3(1-p)^9.
\]
For $D_1$ ($p=\tfrac56$): $L(\tfrac56)\approx 5.74\times10^{-8}$. For $D_2$
($p=\tfrac16$): $L(\tfrac16)\approx 8.97\times10^{-4}$. The data are vastly more
likely under $D_2$, so it is almost surely the mostly-white die.
\end{example}

\section{The Likelihood Function}

For a sample $x_1,\dots,x_n$ from a distribution with parameter $\theta$, the
\emph{likelihood function} is the probability (discrete case) or density
(continuous case) of the data, viewed as a function of $\theta$:
\[
  L(\theta)=\prod_{i=1}^n p_\theta(x_i)\quad\text{(discrete)},\qquad
  L(\theta)=\prod_{i=1}^n f_\theta(x_i)\quad\text{(continuous)}.
\]
The \emph{maximum likelihood estimate} (MLE) is the maximizer
$\hat\theta=\argmax_\theta L(\theta)$.

\begin{example}[The exponential rate]
For a sample from $\Exp(\lambda)$, find the MLE of $\lambda$.

\solution The likelihood is
\[
  L(\lambda)=\prod_{i=1}^n\lambda e^{-\lambda x_i}=\lambda^n e^{-\lambda\sum x_i}.
\]
Differentiating and setting $L'(\lambda)=0$ gives $1-\lambda\xbar_n=0$, so the
MLE is $\hat\lambda=1/\Xbar_n$. (One checks the second derivative is negative
there, confirming a maximum.)
\end{example}

\begin{example}[An endpoint that defies calculus]
A sample $x_1,x_2,x_3=0.98,1.57,0.31$ comes from $\Unif(0,\theta)$ with $\theta$
unknown. Find the MLE of $\theta$.

\solution Each density is $1/\theta$ on $[0,\theta]$ and $0$ outside, so the
likelihood is
\[
  L(\theta)=\begin{cases} 1/\theta^3, & \theta\ge\max(x_i)=1.57,\\ 0, & \theta<1.57.\end{cases}
\]
Differentiating would only show $L$ is decreasing where it is positive; the
maximum sits at the \emph{boundary}. As Figure~\ref{fig:likuniform} shows,
$L(\theta)$ jumps up to its largest value exactly at $\theta=1.57$ and decays
thereafter, so $\hat\theta=\max(x_1,x_2,x_3)=1.57$. In general the MLE for the
upper endpoint is $\max\{X_1,\dots,X_n\}$.
\end{example}

\begin{figure}[ht]
  \centering
  \includegraphics[width=0.5\linewidth]{likelihood_uniform}
  \caption{Likelihood for a $\Unif(0,\theta)$ sample. It is zero until $\theta$
  reaches the largest observation, then decreases—so the maximum is at the
  boundary, where calculus alone would not find it.}
  \label{fig:likuniform}
\end{figure}

\section{The Loglikelihood}

Because $L$ is a product, differentiating it directly means repeated use of the
product rule. Taking logarithms converts the product into a sum, which is far
easier to differentiate; and since the logarithm is increasing, the
\emph{loglikelihood} $\ell(\theta)=\ln L(\theta)$ has its maximum at the same
$\theta$ as $L$.

\begin{example}[Cycles to pregnancy, with censoring]
For $100$ smoking women, the number of cycles to pregnancy is modeled as
$\Geo(p)$; $7$ women had an unknown number exceeding $12$. Counting
$29$ ones, $16$ twos, \dots, and $7$ values above $12$, the likelihood works out
(up to a constant) to
\[
  L(p)=C\,p^{93}(1-p)^{322}.
\]
Find $\hat p$.

\solution The constant $C$ counts arrangements and does not affect the
maximizer. Differentiating,
\[
  L'(p)=C\,p^{92}(1-p)^{321}\bigl(93(1-p)-322p\bigr)=C\,p^{92}(1-p)^{321}(93-415p),
\]
which vanishes (in the interior) at $p=93/415\approx 0.224$. The loglikelihood
$\ell(p)=\ln C+93\ln p+322\ln(1-p)$ has the same maximizer
(Figure~\ref{fig:liksmokers}). Notably $0.224$ is well below the naive estimate
$29/100=0.29$ that ignores the censored women. For the $474$ nonsmokers the
likelihood is $\propto p^{474}(1-p)^{955}$, giving $\hat p=474/1429\approx 0.33$.
\end{example}

\begin{figure}[ht]
  \centering
  \includegraphics[width=\linewidth]{likelihood_smokers}
  \caption{Likelihood $L(p)\propto p^{93}(1-p)^{322}$ (left) and its
  loglikelihood (right) for the smokers. Both peak at the same value
  $\hat p=93/415\approx 0.224$; the loglikelihood is easier to differentiate.}
  \label{fig:liksmokers}
\end{figure}

\section{The Normal and Poisson MLEs}

\begin{example}[Estimating the mean and standard deviation of a normal]
For an i.i.d.\ $\N(\mu,\sigma^2)$ sample, find the MLEs of $\mu$ and $\sigma$.

\solution The loglikelihood is
\[
  \ell(\mu,\sigma)=-n\ln\sigma-n\ln\sqrt{2\pi}-\frac{1}{2\sigma^2}\sum_{i=1}^n(x_i-\mu)^2.
\]
Setting $\partial\ell/\partial\mu=0$ gives $\hat\mu=\xbar_n$, and then
$\partial\ell/\partial\sigma=0$ gives
\[
  \hat\sigma=\sqrt{\frac1n\sum_{i=1}^n(x_i-\xbar_n)^2}.
\]
(The same two answers follow if one parameter is estimated with the other held
fixed—for instance, with $\mu$ known to be $0$ the MLE of $\sigma$ is
$\sqrt{\tfrac1n\sum x_i^2}$.) Notice the MLE of the variance divides by $n$, not
$n-1$; we return to this below.
\end{example}

\begin{example}[Poisson arrivals, and the probability of none]
Packets arrive at a server, modeled as $\Pois(\mu)$ per minute. (a)~Find the MLE
of $\mu$. (b)~Find the MLE of the probability $e^{-\mu}$ of an idle minute.

\solution (a)~The likelihood is
$L(\mu)=\prod_i \tfrac{\mu^{x_i}}{x_i!}e^{-\mu}=\tfrac{e^{-n\mu}\mu^{\sum x_i}}{\prod x_i!}$,
with loglikelihood $\ell(\mu)=(\sum x_i)\ln\mu - n\mu - \ln\prod x_i!$. Then
$\ell'(\mu)=\tfrac{\sum x_i}{\mu}-n=0$ gives $\hat\mu=\xbar_n$, a maximum since
$\ell''<0$. (Applied to the South London bomb data of
Chapter~\ref{ch:estimation}, this gives $\hat\mu=537/576\approx 0.93$, confirming
the earlier estimate as the MLE.) (b)~By the invariance property below, the MLE
of $e^{-\mu}$ is simply $e^{-\xbar_n}$.
\end{example}

\section{Estimating a Population Size}

\begin{example}[The marked-bolt problem, revisited]
Recall the box of $N$ bolts from Chapter~\ref{ch:discrete}. (a)~If draws
\emph{with} replacement give a $\Geo(1/N)$ sample, find the MLE of $N$. (b)~If
draws \emph{without} replacement give a discrete uniform sample on
$\{1,\dots,N\}$, find the MLE of $N$.

\solution (a)~The loglikelihood for $\Geo(1/N)$ data is
$\ell(N)=-n\ln N+\bigl(\sum x_i-n\bigr)\ln(1-\tfrac1N)$; differentiating and
solving $\ell'(N)=0$ yields, after simplification, $\hat N=\xbar_n$. (b)~For the
discrete uniform, $L(N)=(1/N)^n$ when $N\ge\max x_i$ and $0$ otherwise, a
decreasing function of $N$ on its support. The maximum is again at the boundary:
$\hat N=\max\{X_1,\dots,X_n\}$, the largest value seen—just as for the continuous
uniform endpoint.
\end{example}

\section{Properties of Maximum Likelihood Estimators}

Beyond being a uniform recipe, MLEs enjoy several strong properties.

\paragraph{Invariance.} If $\hat\theta$ is the MLE of $\theta$ and $g$ is an
invertible function, then $g(\hat\theta)$ is the MLE of $g(\theta)$. (This is why
$e^{-\xbar_n}$ is the MLE of $e^{-\mu}$, and why the MLE of $\sigma^2$ is the
square of the MLE of $\sigma$.)

\paragraph{Asymptotic unbiasedness.} An MLE may be biased in finite samples. The
normal-variance MLE $D_n^2=\tfrac1n\sum(X_i-\Xbar_n)^2$ equals
$\tfrac{n-1}{n}S_n^2$, so $\E[D_n^2]=\tfrac{n-1}{n}\sigma^2$—biased downward, but
with the bias vanishing as $n\to\infty$. Under mild conditions every MLE is
asymptotically unbiased: $\lim_{n\to\infty}\E[T_n]=\theta$.

\paragraph{Asymptotic minimum variance.} The variance of any unbiased estimator
is bounded below by the \emph{Cram\'er--Rao} bound, and MLEs attain this bound
asymptotically: no estimator does better in large samples.

\paragraph{Asymptotic normality.} The standardized MLE
$(\hat\theta-\theta)/\widehat{\mathrm{se}}$ converges to $\N(0,1)$. Consequently
\[
  \hat\theta_n\pm 1.96\,\widehat{\mathrm{se}}
\]
is an approximate $95\%$ confidence interval for $\theta$.

\begin{example}[Confidence intervals for the pregnancy parameters]
For the $\Geo(p)$ model, the MLE $\hat p$ has standard error
$\mathrm{se}=\sqrt{p^2(1-p)/n}$. Using asymptotic normality, build approximate
$95\%$ intervals for the smokers ($\hat p=0.224$, $n=100$) and nonsmokers
($\hat p=0.33$, $n=486$).

\solution The interval is $\hat p\pm 1.96\sqrt{\hat p^2(1-\hat p)/n}$. For the
smokers this gives roughly $(0.19,0.26)$, and for the nonsmokers roughly
$(0.31,0.35)$. The two intervals do not overlap, strong evidence that smokers and
nonsmokers really do have different per-cycle conception probabilities.
\end{example}
